\section{Counterintuitive Patterns}

Several relationships in the dataset deviate from textbook expectations during
specific periods. Below we identify each pattern, the dates at which it
manifests, and the economic rationale.

\subsection*{CDX IG and UST10: positive correlation}

Under normal conditions, investment-grade credit spreads widen when Treasury
yields fall (flight-to-quality), producing a negative correlation. However, the
rolling 60-day correlation turned positive in two distinct episodes:

\begin{itemize}
  \item \textbf{Q1--Q2 2021} (March--July, peak \(+0.34\)). Vaccine
    deployment and large-scale fiscal stimulus supported a broad reflation in
    asset prices. Yields rose on higher inflation expectations, while IG
    spreads widened as the market repriced the term premium and the likelihood
    of an earlier-than-anticipated start to monetary tightening. Both moved
    higher simultaneously, producing a positive correlation.

  \item \textbf{H2 2022 -- Q1 2023} (June 2022--March 2023, peak \(+0.48\)).
    The Federal Reserve was conducting cumulative rate increases of over
    400\,bps during this period, and both yields and spreads were driven
    higher by the same factor: persistently above-target inflation requiring
    sustained monetary tightening. Higher policy rates raised corporate funding
    costs and recession probability simultaneously, causing credit and rates to
    sell off together---a regime in which the traditional negative spread--yield
    relationship breaks down.
\end{itemize}

\subsection*{Crude oil and CDX HY: sign reversal}

Oil and high-yield credit are typically positively correlated through the growth
channel: robust demand supports both oil prices and HY issuers (a significant
share of which, in the US, are energy producers). The rolling correlation was
positive for most of the sample but reversed in early 2022:

\begin{itemize}
  \item \textbf{Pre-2022}: the correlation was negative to near-zero (averaging
    roughly \(-0.20\) to \(-0.40\)), consistent with an environment where oil
    supply shocks and OPEC dynamics dominated---rising oil prices were
    associated with cost-push inflation fears rather than demand strength,
    weighing on HY credit.

  \item \textbf{March--May 2022} (peak \(+0.46\)). Following Russia's invasion
    of Ukraine, oil prices rose above \$100 on supply disruption concerns. HY
    spreads widened concurrently as recession probability increased. As both
    subsequently declined---oil on weakening demand expectations, HY spreads on
    partial relief from peak tightening fears---the correlation turned strongly
    positive. During this period, oil functioned primarily as a geopolitical
    supply-risk indicator rather than a demand proxy, aligning its
    directional moves with those of credit spreads.
\end{itemize}

\subsection*{S\&P 500 and UST10: persistent sign flips}

Equities and Treasury yields can be positively or negatively correlated
depending on whether the dominant macro driver is growth or inflation:

\begin{itemize}
  \item \textbf{2019 -- mid-2020} (peak \(+0.79\)). The correlation was
    strongly positive as both equities and yields were driven by growth
    expectations. In 2019 the Federal Reserve was reducing rates in response
    to slowing growth and trade-policy uncertainty; equities and yields
    declined together. During the COVID-19 sell-off (February--March 2020) and
    the subsequent recovery, both again moved in tandem---lower on the demand
    shock, higher on the monetary and fiscal policy response.

  \item \textbf{Q1--Q2 2021} (trough \(-0.38\)). The correlation turned
    negative as inflation expectations diverged from equity sentiment. Yields
    rose on reflation repricing while equities were volatile on concerns that
    higher rates would compress valuations---an environment in which
    improving macroeconomic data paradoxically weighed on duration-sensitive
    equities through the discount-rate channel.

  \item \textbf{H2 2022 -- Q1 2023} (trough \(-0.48\)). The correlation was
    persistently negative throughout the period of sustained monetary
    tightening. Higher yields reflected restrictive policy, which directly
    impaired equity valuations. Hawkish policy communications pushed yields
    higher and equities lower, while softer inflation data produced the
    opposite effect. This pattern is characteristic of an inflation-dominated
    regime in which the equity--rate relationship inverts relative to
    growth-dominated periods.
\end{itemize}
